\section{Measurable realizations}
\label{sec:measurable-realizations}

We give the parameterized Dirichlet construction and the measurable version of the bounded phase used above.
The common probability space makes the coefficients jointly measurable in the mass, the base probability, and the distinguished rate.

\subsection{Dirichlet probabilities with varying parameters}

We use the stick-breaking formula of Lemma~\ref{lem:stick-breaking}, choosing the locations by quantiles and the break fractions by common uniform random variables.

\begin{lemma}[A parameterized Dirichlet realization]
\label{lem:parameterized-dirichlet}
There is a fixed probability space \((\Omega,\mathcal A,\mathbb P)\) and jointly Borel maps
\[
 (B,F,\omega)\longmapsto Q(B,F,\omega),\qquad
 (B,F,y,\omega)\longmapsto\widehat P(B,F,y,\omega)
\]
into \(\mathcal P(\R)\), for \(B>0\), \(F\in\mathcal P(\R)\), and \(y\in\R\), such that
\[
 Q(B,F,\cdot)\sim\DP(BF),\qquad
 \widehat P(B,F,y,\cdot)\sim\DP(BF+\delta_y).
\]
The construction uses a single event of probability one on which the stick-breaking weights sum to one for every \(B>0\) and every \(F\).
All spaces of probabilities carry their narrow Borel sigma-fields.
\end{lemma}

\begin{proof}
For \(F\in\mathcal P(\R)\) and \(0<u<1\), define its quantile by
\[
 q_F(u)=\inf\{x\in\R:F((-\infty,x])\ge u\}.
\]
This is a finite real number.
For every \(a\in\R\),
\begin{equation}\label{eq:quantile-joint-borel}
 \{(F,u):q_F(u)<a\}
 =\bigcup_{\substack{r\in\mathbb Q\\r<a}}
       \{(F,u):F((-\infty,r])\ge u\}.
\end{equation}
The evaluation \(F\mapsto F((-\infty,r])\) is Borel: the bounded continuous functions \((1-n(x-r)_+)_+\) decrease to the indicator of this half-line, so their integrals converge to that evaluation.
Thus \((F,u)\mapsto q_F(u)\) is jointly Borel.
Right continuity of distribution functions also gives \(\mathbb P\{q_F(S)\le x\}=F((-\infty,x])\) for a uniform \(S\) on \((0,1)\).

Take \(\Omega=(0,1)^{\mathbb N}\times(0,1)^{\mathbb N}\times(0,1)\) with its product Borel sigma-field and product uniform probability.
Write \((S_j)_j,(T_j)_j,T_*\) for its independent coordinate variables.
Suppressing \(\omega\) in the formulas, put
\begin{equation}\label{eq:parameterized-stick-breaking}
 \begin{gathered}
 Z_j(F)=q_F(S_j),\qquad V_j(B)=1-T_j^{1/B},\\
 W_j(B)=V_j(B)\prod_{i<j}(1-V_i(B)).
 \end{gathered}
\end{equation}
These are jointly Borel in their parameters and \(\omega\).
The decreasing products \(\prod_{j\le m}T_j\) tend to zero almost surely: their expectations are \(2^{-m}\), so their nonnegative limit has expectation zero.
Denote this Borel event by \(\Omega_0\).
It does not depend on \(B,F\), and on it
\[
 R_m(B):=\prod_{j\le m}(1-V_j(B))
        =\left(\prod_{j\le m}T_j\right)^{1/B}
        \longrightarrow0\qquad\text{for every }B>0.
\]
In particular \(\sum_jW_j(B)=1\) simultaneously for all parameters on \(\Omega_0\); no intersection of parameter-dependent null sets is needed.

The probability-valued maps
\[
 Q_m(B,F)=\sum_{j\le m}W_j(B)\delta_{Z_j(F)}+R_m(B)\delta_0
\]
are jointly Borel, since finite convex combinations and \(z\mapsto\delta_z\) are continuous in the narrow topology.
On \(\Omega_0\) they converge narrowly to \(Q(B,F)=\sum_{j\ge1}W_j(B)\delta_{Z_j(F)}\); the tail error against a bounded continuous test is at most twice its supremum norm times \(R_m(B)\).
Set both \(Q_m\) and \(Q\) equal to \(\delta_0\) on \(\Omega_0^c\).
The resulting everywhere pointwise limit is a jointly Borel map into the metrizable space \(\mathcal P(\R)\).
For each fixed \(B,F\), its law is \(\DP(BF)\) by Lemma~\ref{lem:stick-breaking}.

Finally, set \(Z_*(B)=1-T_*^{1/B}\) and
\[
 \widehat P(B,F,y)=(1-Z_*(B))Q(B,F)+Z_*(B)\delta_y.
\]
This map is jointly Borel.
For each fixed \(B,F\), the variable \(Z_*(B)\) is \(\operatorname{Beta}(1,B)\) and independent of \(Q(B,F)\).
Lemma~\ref{lem:dirichlet-posterior}, in log-rate coordinates, therefore gives the stated posterior law.
Passing through \(\exp_*\) supplies the corresponding jointly Borel realizations on positive rates as well.
\end{proof}


\subsection{Proof of Lemma~\ref{lem:phase}}
\label{sec:phase-proof}

\begin{proof}[Proof of Lemma~\ref{lem:phase}]
The function \(M_P\) is a nonzero Stieltjes function, since \(\int(1+b)^{-1}P(\dd b)\le1\).
For \(z\) in the upper half-plane,
\[
  \operatorname{Im}M_P(z)
  =-\operatorname{Im}z\int|z+b|^{-2}P(\dd b)<0.
\]
Thus the arguments in \eqref{eq:phase-representative} are unambiguous principal arguments, with \(-\pi<\arg M_P(z)<0\).
Apply \eqref{eq:ssv-exponential} to \(f=1/M_P\), and subtract its real logarithms at \(s\) and \(1\).
Negating the resulting identity gives exactly the sign in \eqref{eq:phase-anchor}.
The uniqueness assertion follows from uniqueness in \eqref{eq:ssv-exponential}: the value at \(1\) fixes the otherwise free multiplicative constant.

Let \(\xi_P^0\) be any almost-everywhere phase initially provided by that representation.
Analytic continuation of \eqref{eq:ssv-exponential} and its imaginary part give
\[
  \arg(1/M_P(-t+i\varepsilon))
  =\int_0^\infty
       \frac{\varepsilon\,\xi_P^0(u)}
            {(u-t)^2+\varepsilon^2}\,\dd u,
  \qquad \varepsilon>0.
\]
Extend \(\xi_P^0\) by zero to the negative half-line.
The normalized integral is its Poisson approximate identity; at each Lebesgue point it tends to \(\xi_P^0(t)\).
For completeness, this convergence follows by splitting the integral into a neighborhood of the point, where the Lebesgue-point averages of the error tend to zero, and its complement, whose Poisson mass tends to zero.
Boundedness of \(\xi_P^0\) controls the complement.
Hence \eqref{eq:phase-representative} agrees with \(\xi_P^0\) almost everywhere and preserves \eqref{eq:phase-anchor}.

For fixed \(n\), the map \((P,t)\mapsto M_P(-t+i/n)\) is jointly Borel (in fact continuous).
The kernel is bounded by \(n\), is continuous in the rate variable, and its dependence on \(t\) is locally uniform, with derivative bounded by \(n^2\).
Taking a continuous argument on the strict lower half-plane and then a limsup proves joint Borel measurability of the specified representative.
The minus sign remains inside the limsup in \eqref{eq:phase-representative}, so the definition is unambiguous also at exceptional boundary points.

Finally,
\[
  \int_0^\infty
   \left|\frac1{s+t}-\frac1{1+t}\right|\dd t=|\log s|,
\]
and, for every integer \(k\ge1\),
\[
  \partial_s^k\log M_P(s)
  =(-1)^k k!\int_0^\infty
                    \frac{\xi_P(t)}{(s+t)^{k+1}}\,\dd t,
  \qquad
  k!\int_0^\infty(s+t)^{-k-1}\dd t
       =(k-1)!s^{-k}.
\]
These estimates prove the absolute convergence and justify all the stated differentiations.
\end{proof}
